
\appendix

%%==================================================================
\section{Proof of the contraction theorem}\label{app:contraction}
%%==================================================================

This appendix gives the proof of Theorem~\ref{thm:contraction}. The argument
and its parameters are those of \oai{} \cite[Appendix D.4]{OAI26}; the
presentation follows theirs step by step, with some computations written out
in more detail. Throughout, $\log$ is to base two and $\ln$ is natural. A call
of the recursion receives a bipartite network with $M$ edges and a working
precision $p$, and returns its contraction modulo $2^p$. One side of the
network is called the row side and the other the column side; the ports of
row vertices are \emph{row ports} and those of column vertices \emph{column
ports}.

\subsection{The coefficient ring}\label{app:ring}
Put $\lambda=\lceil\log s\rceil$ and $e=4\lambda$. Choose a monic $f\in\Z[U]$ of
degree $e$ with coefficients in $\{0,1\}$ whose reduction modulo two is
irreducible; such $f$ exists because $\F_{2^e}$ exists, and it is found in time
polynomial in $s$ by testing the $2^e=s^{O(1)}$ candidates for divisibility by
every monic polynomial of degree at most $e/2$ over $\F_2$ \cite{vzGG13}. Put
\[
\cO=\Z_{(2)}[U]/(f),\qquad K=\Q[U]/(f),\qquad \cA=\cO[V]/(V^\ell),\qquad
\cA_p=\cA/(2^p),
\]
where $\Z_{(2)}$ is the ring of rationals with odd denominator. Since $f$ is
irreducible modulo two it is irreducible over $\Q$, so $K$ is a field. For
$x\in K^\times$ let $\vtwo(x)$ be the least two-adic valuation of its coordinates
in the basis $1,U,\dots,U^{e-1}$. Then $\cO/2\cO=\F_{2^e}$, and factoring out the
least coordinate valuation shows $\vtwo(xy)=\vtwo(x)+\vtwo(y)$. An element of
$\cO$ that is nonzero modulo two is a unit, since the determinant of its
multiplication matrix is odd. Arithmetic in $\cA_p$, with coordinates in the
bases $1,U,\dots,U^{e-1}$ and $1,V,\dots,V^{\ell-1}$ stored as residues in
$[0,2^p)$, takes time and space polynomial in $s,p,\ell$; units are inverted by
inversion modulo two and Newton lifting. Divisibility in $\cA$ is coordinatewise.

The lifts to $\cO$ of the $2^e\ge s^4$ elements of $\F_{2^e}$, taken with
coordinates in $\{0,1\}$, have pairwise differences that are units; Lagrange
interpolation at fewer than $s^4$ of them is therefore valid over every
$\cA_p$. The original entries lie in $B_{P_0,\ell}\subseteq\cA_{P_0}$, and the
contraction, a polynomial in the entries, is computed in $\cA_{P_0}$ and read off
in $B_{P_0,\ell}$.

When the precision grows, each tensor entry is lifted separately: an entry
specified modulo $2^{p_c}$ is replaced by the element of $\cA$ whose coordinates
are its integer representatives in $[0,2^{p_c})$. This lift is not additive, but
it is fixed: every tensor carries the precision $p_c$ at which it was created,
and every later evaluation uses the same representatives. Precisions strictly
increase along a recursive path, so the representatives are canonical at every
later precision.

\subsection{Parameters and packing}\label{app:packing}
All inequalities below hold for $s$ larger than a constant depending on
$C,A,b$; smaller sizes are handled by direct summation. Put
\[
H=\Bigl\lceil\frac{s}{\lfloor\sqrt\lambda\rfloor}\Bigr\rceil,\qquad r=64,\qquad
\varepsilon=\frac1{128\cdot1024},\qquad\Delta=\lceil s^{9/10}\rceil.
\]
A call with $M\le H$ edges is a \emph{leaf} and is evaluated by summing over
all $2^M$ edge assignments. At an \emph{active} call at depth $d$ (the root has
depth zero) put
\begin{equation}\label{eq:params}
u=u_d=\Bigl\lfloor\frac{\varepsilon\lambda}{2^d}\Bigr\rfloor,\qquad
a=a_d=2\sum_{j=0}^du_j,\qquad L=2^a,\qquad \Darr=r^{2a}.
\end{equation}
It will be shown that every child of a call with $M$ edges has at most $M/8$
edges and maximum degree at most $\Delta$. Hence on every active path
$H<M\le Cs8^{-d}$, so $8^d<C\lfloor\sqrt\lambda\rfloor+1$ and
$d\le\frac16\log\lambda+O(1)$. Consequently, on every active call,
\begin{equation}\label{eq:inv1}
u=\Theta(\lambda/2^d),\qquad \lambda^{4/5}\le u\le M,\qquad a\le4\varepsilon\lambda,
\qquad L,\Darr\le s^{1/1000},
\end{equation}
since $2^d=O(\lambda^{1/6})$, $a\le4u_0$ and $\Darr=2^{12a}\le2^{48\varepsilon\lambda}$.
In particular
\begin{equation}\label{eq:inv2}
M\ge2\Darr^2\Delta,\qquad M\ge L\Darr,\qquad 2L\le\Delta,
\end{equation}
because $M>H\ge s^{1-o(1)}$ while the right sides are $O(s^{9/10+1/500})$.

\begin{lemma}[Packing]\label{lem:packing}
Assume \eqref{eq:inv2}. One can label the vertices on each side by elements of
$[r]^a$ and place the edges into \emph{packets} of $L$ slots so that:
\begin{enumerate}[label=\textup{(\arabic*)},leftmargin=2.2em]
\item there are at most $3M/L$ packets, every edge occupies the same slot at its
row end and at its column end, and unused slots carry no ports;
\item in every packet and every coordinate $c\in[a]$ the unordered pair of row
and column symbols is fixed, and when the two symbols differ, bit $c$ of the
slot number records which of them is the row symbol;
\item for every coordinate and every symbol, at most $3M/r$ real ports on each
side belong to vertices with that symbol in that coordinate.
\end{enumerate}
The construction takes time and space polynomial in $s+L$.
\end{lemma}

\begin{proof}
The \emph{pair label} of an edge is the pair (row label, column label), one of
$\Darr$ values; let $X_\gamma$ be the number of edges with pair label $\gamma$.
For independent uniform labels,
\[
\mathbb E\sum_\gamma\Bigl(X_\gamma-\frac M{\Darr}\Bigr)^2
=\sum_{e,e'}\Bigl(\Pr[\gamma(e)=\gamma(e')]-\frac1{\Darr}\Bigr),
\]
the sum running over ordered pairs of edges. If $e$ and $e'$ have disjoint
endpoints their labels are independent and the term vanishes. At most $2M\Delta$
ordered pairs share an endpoint, and each term is at most one, so the
expectation is at most $2M\Delta$. Labelling the vertices one at a time and
choosing, among the $r^a$ labels of the next vertex, one that does not increase
the conditional expectation \cite{Rag88} gives a deterministic labelling with
the same bound; each conditional expectation is a sum over pairs of edges of
probabilities involving at most four vertices, so it is computed exactly in
polynomial time. Every load is then at most $M/\Darr+\sqrt{2M\Delta}$. Pad every
pair-label class to $H_{\mathrm{pad}}$ slots, the largest load rounded up to a
multiple of $L$; by \eqref{eq:inv2},
\[
H_{\mathrm{pad}}\le\frac M{\Darr}+\sqrt{2M\Delta}+L\le\frac{3M}{\Darr}.
\]
Group the pair labels by the unordered pair of symbols in each coordinate. If
$a'$ coordinates of a group have unequal symbols, the group consists of $2^{a'}$
pair labels, which differ in orientation. Form packets of $L$ slots by taking
$2^{a-a'}$ slots of each of these $2^{a'}$ labels; this is possible since
$H_{\mathrm{pad}}$ is divisible by $L$. Number the positions of a packet from $0$
to $L-1$ so that bit $c$ gives the orientation when the pair in coordinate $c$ is
unequal and indexes copies otherwise, with one convention throughout a packet.
There are $H_{\mathrm{pad}}\Darr/L\le3M/L$ packets. In a fixed coordinate a fixed
symbol is the row symbol of exactly $1/r$ of the padded pair labels, so it has at
most $H_{\mathrm{pad}}\Darr/r\le3M/r$ real row ports, and likewise for columns.
\end{proof}

\subsection{A Cauchy deformation}\label{app:cauchy}
Fix an active call, its packets and its lifted tensors. Put
\[
S=\{-L,-L+1,\dots,L-1\},\qquad C(T)=\prod_{t\in S}(T-t),\qquad \alpha=2^{a+1}.
\]
For a row slot $x$ and a column slot $y$ of the same packet put
\begin{equation}\label{eq:W}
W_{xy}(T)=\frac{C(T)}{C'(0)\,(T-(y-x))},
\end{equation}
a polynomial of degree $2L-1$ over $\Q$, since $y-x\in S$; entries between
different packets are zero, and $W$ is restricted to real ports. At $T=0$ the
diagonal entries are one and the others vanish, because $C(T)/T\to C'(0)$ and
every other quotient keeps the factor $T$; so $W(0)$ is the identification of
row and column ports given by the edges. Let $A(Y)$ be the product of the
lifted row tensors on a row assignment $Y\in\{0,1\}^M$ and $B(X)$ the product of
the lifted column tensors on a column assignment $X$. With $l_0(z)=1-z$ and
$l_1(z)=z$ put
\begin{equation}\label{eq:F}
F(T)=\sum_{X,Y\in\{0,1\}^M}A(Y)B(X)\prod_{\text{row ports }x}l_{Y_x}\bigl((W(T)X)_x\bigr).
\end{equation}
Its coefficients lie in $K[V]/(V^\ell)$, its degree is at most $M(2L-1)$, and,
since $\prod_xl_{Y_x}(X_x)=[Y=X]$, $F(0)=\sum_XA(X)B(X)$ is the contraction of the
lifted tensors.

\begin{lemma}[Integrality]\label{lem:integral}
For every $t\in S$ the entries of $W(t+\alpha z)$ lie in $\cO[z]$ and
$F(t+\alpha z)\in\cA[z]$. Given the coordinates of $\zeta$ modulo $2^p$, each
$W_{xy}(t+\alpha\zeta)\bmod2^p$ is computed in time polynomial in $s,p$.
\end{lemma}

\begin{proof}
$S$ is a complete residue system modulo $2^{a+1}$. For $q\ne t$ in $S$ one has
$\vtwo(t-q)\le a$, and the number of $q$ with $\vtwo(t-q)=k$ does not depend on
$t$; hence $\vtwo(C'(t))=\vtwo(C'(0))$. In
\[
\frac{C(t+\alpha z)}{C'(0)}=\frac{\alpha z\prod_{q\ne t}(t-q+\alpha z)}{C'(0)}
\]
each factor $t-q+\alpha z$ is divisible coefficientwise by $2^{\vtwo(t-q)}$, so
after cancelling these powers against the denominator the polynomial is
divisible by $2^{a+1}$. Formula~\eqref{eq:W} omits the factor indexed by
$q=t-(y-x)$. If $y-x=t$ the omitted factor is $\alpha z$ and the rest is
integral; otherwise $\alpha z$ remains, the omitted factor has content
$2^{\vtwo(t-y+x)}$ with $\vtwo(t-y+x)\le a$, and the quotient is still
integral. Substituting in \eqref{eq:F} gives the claim for $F$. To evaluate,
divide each factor $t-q+\alpha z$ by $2^{\vtwo(t-q)}$, cancel the remaining
powers of two against $C'(0)$, invert its odd part modulo $2^p$, and multiply
the $O(L)$ factors at $\zeta$.
\end{proof}

\subsection{Tapered Hermite recovery}\label{app:hermite}
\begin{lemma}[Recovery]\label{lem:recovery}
Let $1\le u\le a$ and $M\ge u$, and put $h=\lceil3M/u\rceil$. For $t\in S\setminus\{0\}$
let
\[
i=\vtwo(t),\qquad q_t=q(i)=h\max(0,u-i),\qquad D_t=L/2^i .
\]
Let $F\in(K[V]/(V^\ell))[T]$ have degree at most $M(2L-1)$, with $F(0)\in\cA$ and
$F(t+\alpha z)\in\cA[z]$ for all $t\in S$. Then $F(0)\bmod2^p$ is determined by,
and computable in polynomial time from, the first $q_t$ coefficients of
\[
F(t+\alpha z)\bmod2^{p_t},\qquad p_t=p+8hD_t+8M(a+1),
\]
for the $t$ with $q_t>0$.
\end{lemma}

\begin{proof}
There are $L$ odd points in $S$, each with $q_t=hu\ge3M$, so
$\sum_{t\ne0}q_t\ge3ML>M(2L-1)$; also $q_t\le hu\le4M$. For $t$ with $q_t>0$ put
\[
E_t(T)=\prod_{q\in S\setminus\{0,t\}}\Bigl(\frac{T-q}{t-q}\Bigr)^{q_q}.
\]
We claim
\begin{equation}\label{eq:hermite}
F(0)=\sum_{t:\,q_t>0}E_t(0)\,\Bigl[\frac{F(t+\alpha z)}{E_t(t+\alpha z)}\Bigr]_{<q_t}\Bigl(-\frac t\alpha\Bigr),
\end{equation}
where $[\cdot]_{<q}$ truncates a power series in $z$ below degree $q$. Replace $0$
by an indeterminate $T$ and $-t/\alpha$ by $(T-t)/\alpha$ on the right. The
summand of $t$ agrees with $F$ to order $q_t$ at $T=t$ and vanishes to order
$q_{t'}$ at every other node $t'$, and its degree is below $\sum q_t$. So the
difference with $F$ is divisible by $\prod_{t\ne0}(T-t)^{q_t}$ and has smaller
degree, hence vanishes; this is Hermite interpolation by the Chinese remainder
theorem \cite[Section 5.6]{vzGG13}.

Each factor of $E_t(t+\alpha z)$ is a power of $1+\alpha z/(t-q)$, which is
integral with constant term one, so the reciprocal truncated below $q_t$ is
integral and the quotient in \eqref{eq:hermite} loses no precision.

Fix $i=\vtwo(t)$ and $D=D_t$. Points $q$ with $\vtwo(q)<i$ have
$\vtwo(t-q)=\vtwo(q)$ and contribute nothing to $\vtwo(E_t(0))$. For $1\le j\le
a-i$ there are $D/2^j$ points with $\vtwo(q)=i+j$, for which $\vtwo(t-q)=i$, and
$D/2^j$ points with $\vtwo(t-q)=i+j$, for which $\vtwo(q)=i$; these exhaust the
points of valuation at least $i$ other than $t$, since the nonzero differences
have valuation at most $a$. Therefore
\begin{equation}\label{eq:vEt}
\vtwo(E_t(0))=\sum_{q\ne0,t}q_q\bigl(\vtwo(q)-\vtwo(t-q)\bigr)
=D\sum_{j=1}^{a-i}\frac j{2^j}\bigl(q(i+j)-q(i)\bigr)\ \ge\ -6Dh,
\end{equation}
because $q(i+j)-q(i)\ge-hj$ and $\sum_{j\ge1}j^2/2^j=6$. Evaluating a polynomial of
degree below $q_t\le4M$ at $-t/\alpha$, whose valuation is at least $-(a+1)$,
loses at most $4M(a+1)$ powers of two. Hence replacing each quotient coefficient
by any lift of its residue modulo $2^{p_t}$ changes the summand by an element of
$2^p\cA$, since $p_t-p>6Dh+4M(a+1)$. The multipliers $E_t(0)(-t/\alpha)^j$ are
rationals of polynomial bit length, so the combination is carried out exactly
with a common denominator, after which the powers of two are cancelled and the
result reduced modulo $2^p$.
\end{proof}

The tapering of the multiplicities is what makes the precision affordable. If
every level below $u$ received the same multiplicity $q_*=\Theta(M)$, the level
$i=u-1$ would lose $q_*D\sum_jj/2^j=\Theta(MD)$ powers of two in
\eqref{eq:vEt}; with multiplicities decreasing by $h=\Theta(M/u)$ per level the
loss is $O(MD/u)$.
Figure~\ref{fig:taper} compares the two choices.

\begin{figure}[t]
\centering
\includegraphics[width=0.8\textwidth]{fig_taper.png}
\caption{Multiplicities of the Hermite nodes by two-adic level, for $u=8$ and
$h=3$. The tapered choice $q(i)=h\max(0,u-i)$ of Lemma~\ref{lem:recovery}
(front) against a flat choice $q_*=hu$ on every level below $u$ (back). Both
give the odd points the multiplicity $hu$ that the degree count needs; only the
tapered one keeps the loss \eqref{eq:vEt} at $O(Dh)$.}\label{fig:taper}
\end{figure}

At a call of precision $p$ apply Lemma~\ref{lem:recovery} with the parameters
\eqref{eq:params}. For each $t$ with $q_t>0$ obtain the whole polynomial
$F(t+\alpha z)\bmod2^{p_t}$ by evaluating it at $M(2L-1)+1<s^4$ of the
interpolation points of Section~\ref{app:ring} and interpolating; the values
$F(t+\alpha\zeta)$ are computed in Section~\ref{app:modes} by child calls of
precision $p_t$. The precisions satisfy
\begin{equation}\label{eq:prec}
p\le s\lambda^{O(1)}2^{a_{d-1}},\qquad p_t=O_C(MD_t/u).
\end{equation}
The first bound holds at the root, where $p=P_0\le Cs$, and the increments
$8hD_t+8M(a+1)$ along a path are at most $s\lambda^{O(1)}2^{a_d}$ at depth $d$,
with $O(\log\lambda)$ depths. For the second, $q_t>0$ forces $i<u$, so
$D_t\ge2^{a_d-u}=2^{a_{d-1}+u}$; since $u\ge\lambda^{4/5}$, the factor $2^u$
exceeds every fixed power of $\lambda$, and both $p$ and $8M(a+1)$ are
$O(MD_t/u)$, as is $8hD_t$. All precisions are polynomial in $s$ because
$a\le4\varepsilon\lambda$.

\subsection{Modes for interactions between groups}\label{app:modes}
Fix $t$ with $i=\vtwo(t)<u$, put $D=L/2^i$, fix an interpolation point $\zeta$, and
work modulo $2^{p_t}$ with $T=t+\alpha\zeta$. Group the vertices on both sides by
their symbol in coordinate $i$; there are $r$ groups. Entries of $W(T)$ whose two
ports lie in the same group are \emph{local}.

The matrix $W(t)$ is local. Indeed $W_{xy}(t)=0$ unless $y=x+t$, and adding $t$,
whose lowest set bit is bit $i$, flips bit $i$ of the slot number. When the pair
of symbols of the packet in coordinate $i$ is unequal, bit $i$ of a slot
records the orientation, so a row slot and a column slot lie in the same group
exactly when their bits $i$ differ; when the pair is equal, the whole packet is
one group.

In a packet with unequal pair, write the column slots as $y=y_0+2^{i+1}j$ with
$0\le y_0<2^{i+1}$; the group of $y$ depends only on $y_0$, so these
\emph{tiles} are monochromatic. If a row $x$ and the tile of $y_0$ lie in
different groups, then $\vtwo(t-y_0+x)\le i$, for otherwise $y_0\equiv x+t\bmod
2^{i+1}$ and their bits $i$ would differ. The same holds with $T$ in place of $t$,
since $\vtwo(\alpha\zeta)>i$. Expanding geometrically,
\begin{equation}\label{eq:modes}
W_{xy}(T)\equiv\sum_{\substack{b\ge0\\ b+1<p_t}}
\frac{C(T)}{C'(0)(T-y_0+x)}\Bigl(\frac{2^{i+1}}{T-y_0+x}\Bigr)^bj^b
\pmod{2^{p_t}}.
\end{equation}
The first factor has valuation at least $a+1-i\ge1$ by the proof of
Lemma~\ref{lem:integral}, and the ratio has valuation at least one, so the term
of index $b$ is divisible by $2^{b+1}$ and the remainder of the geometric series
by $2^{p_t}$. For each packet, tile and $b$ introduce a \emph{mode} $Z$ of weight
$w_Z=b+1$, standing for
\[
L_Z(X)=\sum_{\text{real columns }y\text{ in the tile}}j^bX_y\qquad(j^0=1).
\]
Its \emph{source} is the group of the tile and its \emph{target} the other group
of the packet's pair. By \eqref{eq:modes}, each row form $(W(T)X)_x$ is its
local part plus a combination of the modes targeted at the group of $x$, with
coefficients divisible by $2^{w_Z}$. A packet has $2^{i+1}$ tiles, so there are at
most
\begin{equation}\label{eq:modecount}
\frac{3M}L\,2^{i+1}=\frac{6M}D
\end{equation}
modes of each weight.

Expand the product in \eqref{eq:F} as a polynomial in the modes. The coefficient
of a monomial $Z^\beta$ is divisible by $2^{\sum_Zw_Z\beta_Z}$, so modulo $2^{p_t}$
only the exponent vectors with
\begin{equation}\label{eq:budget}
\sum_Zw_Z\beta_Z<p_t
\end{equation}
contribute. Their number is $\exp(O(\sqrt{Mp_t/D}))=\exp(O_C(M/\sqrt u))$. Indeed,
for $\theta>0$ the generating function bounds the number by
\[
\exp\Bigl(\theta p_t+\frac{6M}D\sum_{j\ge1}-\ln(1-e^{-\theta j})\Bigr),\qquad
\sum_{j\ge1}-\ln(1-e^{-\theta j})=\sum_{h\ge1}\frac1{h(e^{\theta h}-1)}\le\frac{\pi^2}{6\theta},
\]
and $\theta=\sqrt{M/(Dp_t)}$ gives $\exp((1+\pi^2)\sqrt{Mp_t/D})$; then
\eqref{eq:prec} gives $\sqrt{Mp_t/D}=O(M/\sqrt u)$. The vectors are enumerated
with polynomial delay by choosing exponents one variable at a time within the
remaining budget.

\subsection{The child networks}\label{app:children}
For one vector $\beta$ satisfying \eqref{eq:budget} and one group $g\in[r]$, the
child network $\cN_{g,\beta}$ is built as follows. On one side put all current
vertices of group $g$, from both sides, with their ports and their lifted
entries reduced modulo $2^{p_t}$. On the other side put one vertex for each
packet, joined to every port of group $g$ in that packet; its ports are the row
bits $Y$ and column bits $X$ of group $g$ in the packet. Its tensor is the
coefficient of $\prod_ZZ^{\beta_Z}$, over the modes $Z$ of the packet targeted at
$g$, in the product of the row factors $l_{Y_x}$ of group $g$ in the packet,
written in the local column bits and the formal modes, multiplied by
$\prod_ZL_Z(X)^{\beta_Z}$ over the modes of the packet whose source is $g$. Empty
products are one; a positive exponent on a mode with no target rows, or on a
mode whose tile has no real column, gives zero.

Every mode has one source and one target, the coefficient extraction factors
over the groups of the rows, and the substitution of each mode uses only bits of
its source group. Distributing the product in \eqref{eq:F} and summing the bits
of each group independently therefore gives
\begin{equation}\label{eq:children}
F(t+\alpha\zeta)\equiv\sum_{\beta}\ \prod_{g=1}^rZ(\cN_{g,\beta})\pmod{2^{p_t}},
\end{equation}
the sum over the vectors satisfying \eqref{eq:budget}; the omitted monomials
vanish modulo $2^{p_t}$.

By Lemma~\ref{lem:packing}(3) a child has at most $6M/r=3M/32<M/8$ edges. Old
vertices keep their degrees, and a packet vertex has at most $2L\le\Delta$ ports,
so the invariants of Section~\ref{app:packing} hold inductively. A packet entry
is computed by expanding a product of at most $L$ affine factors, each with
polynomially many terms, in time $\exp(O(L\log s))=\exp(O(s^{9/10}))$ and
polynomial space; powers $L_Z(X)^{\beta_Z}$ are computed by repeated squaring. A
new tensor is described by its packet data, $\beta$, $\zeta$ and its creation
precision, and old tensors keep their descriptions, shared by reference, so
entry routines do not compose along the recursion.

\subsection{Proof of Theorem~\ref{thm:contraction}}\label{app:proof}
Embed the entries in $\cA_{P_0}$ and run the recursion. At an active call, for
every $t$ with $q_t>0$ and every interpolation point, compute
\eqref{eq:children} from the child contractions, interpolate $F(t+\alpha z)$,
and apply Lemma~\ref{lem:recovery}; at a leaf sum directly. Correctness follows
by induction from Lemmas~\ref{lem:integral} and~\ref{lem:recovery} and
\eqref{eq:children}.

A call has polynomially many nodes, points and groups, so by
\eqref{eq:budget} an active call at depth $d$ has at most
\[
s^{O(1)}\exp\Bigl(O_C\Bigl(\frac s{\sqrt\lambda}\Bigl(\frac{\sqrt2}8\Bigr)^d\Bigr)\Bigr)
\]
children, using $M\le Cs8^{-d}$ and $u_d=\Theta(\lambda/2^d)$. Over the
$O(\log\lambda)$ depths the number of calls is at most
$\exp(O_C(s/\sqrt\lambda)+O(\lambda\log\lambda))=\exp(O_C(s/\sqrt\lambda))$. A leaf
costs $2^H$ assignments with polynomially many entry evaluations each, that is
$\exp(O_C(H)+O_{C,A}(s^{9/10}))(s+L)^{O(1)}=\exp(O_{C,A}(s/\sqrt\lambda))(s+L)^{O(1)}$.
All other work is polynomial per generated term, which proves the time bound.
For space, traverse the recursion depth first, keeping at each of the
$O(\log\lambda)$ levels the Taylor data of the current call, its indices and its
new tensor descriptions, all of polynomial size by \eqref{eq:prec}. Finally read
off the coordinates of the result in $B_{P_0,\ell}$. \qed

%%==================================================================
\section{Programs that check the finite identities}\label{app:checks}
%%==================================================================

The archive contains a program written in Python with exact integer and
rational arithmetic. It builds the network $\cN(C)$ of
Section~\ref{sec:counting} for random circuits in $\cK(F,\beta)$ with all gate
kinds of Definition~\ref{def:gates}, including big weighted gates split into
blocks, and compares its contraction with exhaustive counting; it checks the
bounds of Lemma~\ref{lem:size} on the edges and degrees, the reduction of
Corollary~\ref{cor:ilp} on random zero-one programs, the identity
\eqref{eq:vEt} and the bound $\vtwo(E_t(0))\ge-6Dh$ for all admissible small
parameters, the integrality of Lemma~\ref{lem:integral}, the Hermite formula
\eqref{eq:hermite} on random polynomials, and the properties (1) to (3) of
Lemma~\ref{lem:packing} on random bipartite graphs with small parameters. For
the degree bound of Lemma~\ref{lem:size} it evaluates the largest possible
degree $\max\bigl(K+b,\lceil s/K\rceil b+1\bigr)$ for every $s$ up to
$3\cdot10^6$: the bound $\lceil s^{9/10}\rceil$ and the inequality $b\le K$ hold
for every $s\ge1416$, which is the constant of that lemma in this range. The
program runs in about twenty seconds and reports $27$ checks passed and none
failed. These checks confirm the finite statements; the proofs in the text do
not depend on them.
